\chapter*{Appendices}
\addcontentsline{toc}{chapter}{Appendices}

\section*{Appendix 1: Complete Benchmark Data}

All benchmarks were executed with $N=1000$ iterations on Windows 11 with Microsoft Visual C++ 2022 compiler using /O2 /arch:AVX2 flags.

\textbf{Statistical Summary:}

\begin{table}[h]
\centering
\caption{Complete Statistical Data (milliseconds)}
\begin{tabular}{l|r|r|r|r}
\hline
\textbf{Algorithm/Operation} & \textbf{Mean} & \textbf{Min} & \textbf{Max} & \textbf{Std Dev} \\ \hline
ECC KeyGen & 97.76 & - & - & - \\
ECC Exchange & 95.62 & - & - & - \\
Kyber-512 KeyGen (SIMD) & 37.75 & - & - & - \\
Kyber-512 Encaps (SIMD) & 54.37 & - & - & - \\
Kyber-512 KeyGen (no SIMD) & 70.54 & - & - & - \\
Kyber-512 Encaps (no SIMD) & 107.15 & - & - & - \\ \hline
\end{tabular}
\end{table}

\textbf{Hardware Configuration:}

\begin{itemize}
\item \textbf{CPU:} x86-64 processor with AVX2 support
\item \textbf{OS:} Windows 11
\item \textbf{Compiler:} MSVC 2022 (\_MSC\_VER 1900+)
\item \textbf{Python:} 3.x with ctypes
\item \textbf{Optimization Flags:} /O2 /arch:AVX2 /fp:fast
\end{itemize}

\section*{Appendix 2: Protocol Wire Format}

\textbf{Packet Structure}

All packets follow the structure:
\begin{verbatim}
[Header Byte] [Length (2 bytes, big-endian)] [Payload]
\end{verbatim}

\textbf{Message Types:}

\begin{table}[h]
\centering
\begin{tabular}{l|l|l}
\hline
\textbf{Header} & \textbf{Value} & \textbf{Description} \\ \hline
HEADER\_HANDSHAKE & 0x01 & Public key transmission \\
HEADER\_MESSAGE & 0x02 & Encrypted application data \\
HEADER\_SYSTEM & 0x03 & Control/status messages \\
HEADER\_KYBER\_CAPSULE & 0x04 & Kyber ciphertext \\ \hline
\end{tabular}
\end{table}

\textbf{Mode Identifiers:}

\begin{itemize}
\item \textbf{MODE\_ECC (0x01):} ECC key exchange mode
\item \textbf{MODE\_KYBER (0x02):} Kyber encapsulation mode
\end{itemize}

\textbf{Handshake Payload Format:}

For HEADER\_HANDSHAKE packets:
\begin{verbatim}
[Mode Byte] [Public Key Data]
\end{verbatim}

\begin{itemize}
\item ECC: 1 byte mode + 33 bytes compressed public key = 34 bytes total
\item Kyber-512: 1 byte mode + 800 bytes public key = 801 bytes total
\item Kyber-768: 1 byte mode + 1184 bytes public key = 1185 bytes total
\item Kyber-1024: 1 byte mode + 1568 bytes public key = 1569 bytes total
\end{itemize}

\section*{Appendix 3: Source Code and Reproducibility}

\textbf{Repository Access:}

The complete source code for this project is available on GitHub:

\begin{center}
\texttt{https://github.com/Aisenesia/pqc-grad-project/tree/SIMD}
\end{center}

The repository includes all implementation files, benchmarking scripts, client-server demonstration code, and the web-based visualization dashboard.

\textbf{Obtaining the Source Code:}

Clone the repository using Git:
\begin{verbatim}
git clone https://github.com/Aisenesia/pqc-grad-project.git
cd pqc-grad-project
\end{verbatim}

Alternatively, download the source as a ZIP archive from the GitHub repository page.

\textbf{Project Structure:}

Start the server (handles TCP port 8000 and WebSocket port 8080):
\begin{verbatim}
python server.py
\end{verbatim}

Start clients in separate terminals:
\begin{verbatim}
python client.py
\end{verbatim}

Access the web dashboard at \texttt{http://localhost:8080/dashboard/}.

\textbf{Compilation:}

Compile the C library with SIMD optimizations:
\begin{verbatim}
cl /O2 /arch:AVX2 /LD lib.c /Fe:lib.dll
\end{verbatim}

For non-SIMD baseline comparison:
\begin{verbatim}
cl /O2 /LD plain_lib.c /Fe:plain_lib.dll
\end{verbatim}

\textbf{Running Benchmarks:}

Execute performance comparison with 1000 iterations:
\begin{verbatim}
python benchmark_comparison.py -n 1000 -o pytest.csv lib.py plain_lib.py
\end{verbatim}

Results are written to CSV files in the \texttt{pytest.csv/} directory.

Server and Client Execution

Start the server (handles TCP port 8000 and WebSocket port 8080):
\begin{verbatim}
python server.py
\end{verbatim}

Start clients in separate terminals:
\begin{verbatim}
python client.py
\end{verbatim}

Access the web dashboard at \texttt{http://localhost:8080/dashboard/}.

Dependencies

\begin{itemize}
\item Python 3.x
\item Microsoft Visual C++ 2022 (for Windows compilation)
\item No external Python packages required (uses standard library only)
\end{itemize}